\section{Proof details}
\label{app:proofs}

\subsection{Differentiation of Bayes risk}

Fix $\alpha_0<1$.  For $\alpha\in[0,\alpha_0]$,
$f_{0,\alpha}\ge(1-\alpha_0)f_0$.  Directly expanding the differentiated
integrand in Proposition~\ref{prop:direction} reduces it to the sum of
$-f_2+f_1^2/f_0$ and
$(t-\delta_\pi)^2\phiN(\,\cdot-t)$.  The first absolute value is bounded by
$f_2+m^2f_0$, and the second by $4m^2\phiN(\,\cdot-t)$.  Both are integrable.
Dominated convergence therefore justifies the right derivative.  The same
argument applies to a direction given by any bounded-support probability
measure after integrating the atom formula in $t$.

\subsection{Fourth derivative used by Simpson's rule}

Let $u(z)=q(t,z)^2$.  With derivatives taken in $z$,
\begin{align*}
u_0&=q^2,\\
u_1&=2qV,\\
u_2&=2(V^2+q\kappa_3),\\
u_3&=2(3V\kappa_3+q\kappa_4),\\
u_4&=2(3\kappa_3^2+4V\kappa_4+q\kappa_5).
\end{align*}
Because $\phiN'=-z\phiN$, repeated product differentiation gives
\[
(u\phiN)^{(4)}
=\phiN\!\left[
u_4-4zu_3+6(z^2-1)u_2+4(3z-z^3)u_1
 +(z^4-6z^2+3)u_0\right].
\]
Every posterior quantity is evaluated on the same $(t,z)$ box.

\subsection{Curvature tail}

For posterior support in $[-m,m]$, $0\le V\le m^2$ by the
Bhatia--Davis inequality.  Also
\[
 |\kappa_3|=\left|\E(\Theta-\mu)^3\right|
 \le 2m\,\E(\Theta-\mu)^2=2mV.
\]
Thus the absolute curvature state is no greater than
\[
C_m=\max\{1,(m^2-1)^2\}+4m^4.
\]
Outside $[-A,A]$, the magnitude omitted from \eqref{eq:curvature} is at most
$2C_m\Pr(|Z|>A)\le4C_m\phiN(A)/A$.

\subsection{Branch pruning invariant}

At every time the active heap plus already pruned cells covers $[0,m]$.
An active cell has a proved upper bound.  A cell is pruned only if its bound
is no larger than the current incumbent plus the work tolerance.  Therefore
the global supremum is bounded by the maximum of the active-cell bounds and
the incumbent plus the work tolerance.  This maximum is exactly the value
reported by the verifier.  The use of floating-point heap priorities affects
only processing order; each retained mathematical bound is an Arb ball.
